% Owner: Afzal. Included from methods.tex. Numbers come from results/tables/E3_*.csv and E4_*.csv.
\subsection{Tasks, folds and features}
\paragraph{T1: next-day ticket counts.} One row per (station, day $d$); the target $y$ is the number of tickets on day $d$ and every feature uses only days before $d$, so each row is a one-step-ahead forecast from history through $d-1$. Features are the counts 1, 7 and 14 days earlier, the means of the previous 7 and 28 days, the mean of all earlier same weekdays, weekday, weekend and Indian-holiday flags of the target day, and the station. The first 28 days have no full history and are dropped, leaving 6{,}642 rows (54 stations, 123 days).

\paragraph{T2: next-week hotspots.} One row per (H3 resolution-8 hexagon, origin week $w$) over the 21 full weeks. The label is 1 if the hexagon is in the top 10\% of hexagons in week $w+1$ by weekly impact, the sum over its tickets of severity times a peak-hour weight (the PCIS trend term is excluded because it is computed over the whole period and would leak the future). Features are the weekly impact of weeks $w$ to $w-3$ and their mean, the evening share, mean severity and junction share of week $w$, and the mean impact of the six neighbouring hexagons.

\paragraph{Folds.} Four expanding-window folds, always training strictly before testing: T1 tests on 14-day blocks starting 13 and 27 February and 12 and 26 March (the last ends on 8 April); T2 tests on 2-week blocks starting at weeks 13, 15, 17 and 19.

\subsection{Models}
\paragraph{T1.} Baselines are \emph{seasonal naive} (the count one week earlier) and a 7-day moving average. Models are a Poisson GLM with log link; ridge \citep{hoerl1970ridge} and lasso \citep{tibshirani1996lasso} on $\log(1+y)$, with count features also log-transformed and the station one-hot encoded; a random forest \citep{breiman2001random} with Poisson splitting criterion, which replicates the base paper \citep{gao2019parking}; and LightGBM \citep{ke2017lightgbm} with the hackathon parameters and a Poisson objective.

\paragraph{T2.} Baselines are \emph{persistence} (last week's label) and \emph{persistence by impact} (percentile rank of last week's impact, the stronger baseline). Models are L2 logistic regression, a linear SVM with sigmoid calibration, a random forest and LightGBM; linear models receive log-transformed impact features.

\paragraph{Tuning.} Each model has one knob (penalty, leaf size or number of trees) with 3 to 5 candidate values, chosen on the last 14 days (T1) or last 2 weeks (T2) of the training window, never on the test block. A first T1 grid with three values left several models on its edge, so the grids were widened once; the chosen values are saved in \texttt{E3\_chosen\_params.csv} and \texttt{E4\_chosen\_params.csv}. The random forest and LightGBM are run with seeds 0, 1 and 2; the other models are deterministic.

\paragraph{Statistics.} T1 uses the macro MASE of the Metrics subsection and paired Wilcoxon tests on per-station MASE; the 12 tests are not corrected for multiple comparisons. T2 uses PR-AUC with a 95\% cluster-bootstrap interval that resamples whole hexagons (1{,}000 draws), F1 at the threshold with the best F1 on the inner validation block, and the Brier score. Every T2 result is shown for all hexagons and for the \emph{emerging} subset, hexagons that were not hotspots in week $w$.
